\section{Discussion}

The central contribution is a boundary on interpretation. A crop ranking can prioritize alternatives without determining the scale at which they should be used. Moving from rank to acreage requires cardinal payoffs, a feasible set, a risk convention and a rule for selecting among multiple optima. This distinction is easy to obscure because both objects can be displayed as ordered crop lists, yet they answer different questions.

The theory sharpens that distinction in three ways. It treats reversal as a property that may concern some, all or one selected optimizer; retains every operational dual term in the marginal condition; and replaces scalar dependence intuition with a restricted distributional order. Feasibility alone can force reversal, while a slack risk constraint can leave an unconstrained optimum unchanged. Consequently, observing discordance is insufficient to attribute it to downside risk or dependence.

The empirical evidence is consistent with the practical relevance of the distinction. State-level ranking and acreage orders often disagree in the admitted complete-case sample, and the pattern survives several definitions and temporal checks. At the same time, the wide spread across definitions and the national null show why one headline prevalence number would be misleading. The analysis documents a phenomenon to explain; it does not recover the explanation.

The failed convergence audit further disciplines interpretation. Numerical reproducibility establishes that code reruns and solvers agree, but it does not guarantee that finite-replication probabilities are precise. Keeping the simulation landscape nonheadline protects the main claim from an attractive but underpowered mechanism story. Future simulation work should increase independent replications under the frozen design or preregister a revised design before promoting prevalence or threshold claims.

Several limitations define the present scope. The state panel covers only \EmpiricalStartYear--\EmpiricalEndYear{} and includes states with complete observations for all three crops. National prices and costs standardize accounting but do not represent local realized margins. Observed acreage embeds unobserved rotations, contracts, expectations, policies and adjustment costs. Farmer-specific risk tolerance and feasible sets are unavailable. The optimizing mechanism is therefore not identified. The linear model omits fixed costs, endogenous prices, spatial interaction and multi-period rotation dynamics. These limits preclude causal, welfare and revealed-optimality interpretations.

The practical implication is therefore procedural. Ranking systems should be accompanied by the decision environment in which their recommendations will be used. At minimum, users need the cardinal score mapping, constraint provenance, uncertainty representation, selection rule and sensitivity to alternative definitions. Where these inputs are unavailable, a ranking-versus-acreage comparison should be labelled descriptive rather than converted into an efficiency or compliance judgment.
